% articulo-16-bibliografia.tex - Bibliografía CSL.
% (generada por pandoc/citeproc)
%
% Este bloque define los entornos y comandos que citeproc emite al volcar
% las referencias. Siempre debe estar presente aunque el documento no tenga
% bibliografía, para evitar errores de «comando no definido» si citeproc
% genera alguna entrada.
%
% CSLReferences recibe dos argumentos posicionales heredados del estándar
% pandoc-citeproc: #1 = flag de sangría francesa (0/1), #2 = separación entre
% entradas (en unidades de \cslentryspacingunit = \parskip).
% La sangría se fija siempre a \parindent (= \ParIndent), independientemente
% de #1, para mantener coherencia con el resto del documento.
%
% \cslhangindent = \parindent: sangría francesa de cada entrada bibliográfica.
% \csllabelwidth = 3em:        ancho de etiqueta para estilos numéricos (APA, etc.)
% \cslentryspacingunit:        unidad de espaciado entre entradas (= \parskip = 0pt
%                              en esta plantilla; citeproc multiplica por #2).
%
% HIPERVÍNCULOS EN CITAS AUTOR-FECHA (link-citations: true)
%
%   Cuando link-citations: true, citeproc emite cada cita en el texto como
%   \citeproc{ref-CLAVE}{texto visible}. El comando \citeproc usa \hyperlink
%   para saltar al ancla correspondiente en la bibliografía.
%
%   \cslbibitem declara ese ancla con \hypertarget{#2}{} antes de cada
%   entrada, donde #2 es el id de la clave bibliográfica (ref-CLAVE).
%   Sin este \hypertarget, \hyperlink no encuentra el destino y el enlace
%   apunta al inicio de la sección en lugar de a la entrada concreta.
%
%   \citeproc requiere que hyperref esté cargado (articulo-09-colores-enlaces.tex).

\usepackage{calc}
\newlength{\cslhangindent}
\setlength{\cslhangindent}{\parindent}
\newlength{\csllabelwidth}
\setlength{\csllabelwidth}{3em}
\newlength{\cslentryspacingunit}
\setlength{\cslentryspacingunit}{\parskip}
\newcommand{\citeproctext}{}
\makeatletter
\newcommand{\cslbibitem}[2][]{%
  \par
  \hypertarget{#2}{}%
  \setlength{\hangindent}{\cslhangindent}%
  \noindent
}
\let\bibitem\cslbibitem
\makeatother
\newenvironment{CSLReferences}[2]{%
  \begingroup
  \setlength{\parindent}{0pt}
  \setlength{\parskip}{#2\cslentryspacingunit}
}{%
  \endgroup
}
\newcommand{\CSLBlock}[1]{#1\hfill\break}
\newcommand{\CSLLeftMargin}[1]{\parbox[t]{\csllabelwidth}{#1}}
\newcommand{\CSLRightInline}[1]{\parbox[t]{\linewidth - \csllabelwidth}{#1}\break}
\newcommand{\CSLIndent}[1]{\hspace{\cslhangindent}#1}
\newcommand{\citeproc}[2]{\hyperlink{#1}{#2}}
